% !TeX root = ../main.tex

\chapter{参考答案}

\section*{1.1 子群与正规子群}

\begin{answers}{1.1}
    \item
    若 \(z\in Z(G)\) 与 \(g\in G\)，则 \(gzg^{-1}=z\in Z(G)\).
    因此 \(Z(G)\trianglelefteq G\).

    \item
    \(H\cap K\) 非空，因为 \(e\in H\cap K\).
    若 \(x,y\in H\cap K\)，则 \(xy^{-1}\in H\) 且 \(xy^{-1}\in K\)，所以 \(xy^{-1}\in H\cap K\).
    由 \autoref{lem:subgroup-test}，\(H\cap K\leq G\).
\end{answers}

\section*{1.2 商群与群同态}

\begin{answers}{1.2}
    \item
    单位元属于像.
    若 \(a=\varphi(x)\) 与 \(b=\varphi(y)\) 属于像，则
    \begin{align*}
        ab^{-1}=\varphi(xy^{-1})
    \end{align*}
    仍属于像.
    由 \autoref{lem:subgroup-test}，\(\operatorname{im}\varphi\leq H\).

    \item
    自然投影把 \(g\in G\) 送到陪集 \(gN\).
    有 \(gN=N\) 当且仅当 \(g\in N\)，因此它的核正是 \(N\).

    \item
    \(\operatorname{sgn}\) 的核由所有偶置换组成，也就是交错群 \(A_n\).
\end{answers}
